\begin{idstatement}{D.1.14}{Proved}\label{st:D.1.14}
Put $e=\widehat\delta$,
$J_2=\left(\begin{smallmatrix}0&1\\-1&0\end{smallmatrix}\right)$,
and let $P_e$ be the integral stabilizer of $e$ for the form
\[
\Omega=\begin{pmatrix}0&0&0&6\\0&0&1&0\\0&-1&0&0\\-6&0&0&0\end{pmatrix}.
\]
Every element of $P_e$ has a unique expression
\[
g(B,x,t)=\begin{pmatrix}1&0&0\\6x&B&0\\t&-x^tJ_2B&1\end{pmatrix},
\quad B\in SL_2(\Z),\quad x\in\Z^2,\quad t\in\Z,
\]
in blocks of sizes $1,2,1$. Let
$Q(x)=x_1+x_2+x_1x_2$, and let
$\alpha:SL_2(\Z)\to\Z/12$ send
$S=\left(\begin{smallmatrix}0&-1\\1&0\end{smallmatrix}\right)$ to $9$
and $T=\left(\begin{smallmatrix}1&1\\0&1\end{smallmatrix}\right)$ to $1$.
Then
\[
\langle A_1,A_2\rangle
=\{g(B,x,t):t+6Q(x)-\alpha(B)=0\pmod{12}\},
\qquad [P_e:\langle A_1,A_2\rangle]=12.
\]
\end{idstatement}
\begin{proof}
Preservation of $e$ and $\Omega$ forces the first row to be
$(1,0,0,0)$ and the last column to be $e$. The middle block has
determinant one. If its first column below the top entry is $v$,
the symplectic equations force the final middle row to be
$-v^tJ_2B/6$. Integrality and unimodularity of $B$ imply
$v\in6\Z^2$, proving the parametrization. Multiplication gives
\[
g(B,x,t)g(C,y,s)=g(BC,x+By,t+s-6x^tJ_2By).             \tag{D6}
\]
Set $a=\left(\begin{smallmatrix}0&1\\-1&-1\end{smallmatrix}\right)$.
Then $A_1=g(a,(1,-1)^t,-2)$,
$A_2=g(S,(0,-1)^t,3)$ and $Sa=T$.
The presentation
$SL_2(\Z)=\langle a,S:a^3=S^4=1,\ [a,S^2]=1\rangle$
follows from the Euclidean algorithm and the real projective-line
action. Modulo $\{\pm I\}$, $a$ has order three and $S$ order
two; $a,a^2$ send positive numbers into the two negative intervals
separated by $-1$, whereas $S$ sends negative numbers to positive
ones. A reduced alternating word cannot act trivially. Lifting
$S^2=-I$ gives the stated presentation. The values
$\alpha(a)=4$, $\alpha(S)=9$ respect the relations and give
$\alpha(T)=1$.

Modulo two, $Q$ is zero only at the origin. Therefore
$Q(Bx)=Q(x)$, and
$Q(x+y)-Q(x)-Q(y)=x^tJ_2y$ modulo two.
Equation (D6) makes
$\chi(g)=t+6Q(x)-\alpha(B)\pmod{12}$ a homomorphism.
It is onto because $\chi(g(I,0,1))=1$, and
$\chi(A_1)=\chi(A_2)=0$ by substitution.

For equality with its kernel, compute
\[
\begin{aligned}
K&=[A_1,A_2^2]=g(I,(0,-1)^t,6),\\
U&=A_2KA_2^{-1}=g(I,(1,0)^t,-6),\\
V&=K^{-1}=g(I,(0,1)^t,-6),\qquad Z=I+12e\gamma.
\end{aligned}
\]
One has $[U,V]=Z^{-1}$ and
\[
U^rV^sZ^n=g(I,(r,s)^t,-6(r+s+rs)+12n).
\]
Thus $U,V,Z$, all in the monodromy group, generate exactly the
intersection of $\ker\chi$ with the kernel of $g\mapsto B$.
The images $a,S$ generate $SL_2(\Z)$.
For an element of $\ker\chi$ choose a monodromy element with the
same middle block. Their quotient lies in that computed
intersection and hence in the monodromy group. This proves equality
and index twelve. The matrices and commutator are checked exactly
by the attached program.
\end{proof}


% D.1.15 from part_D.tex

